\documentclass[11pt]{article}

\usepackage[margin=1in]{geometry}
\usepackage{amsmath,amssymb,amsthm}
\usepackage{graphicx}
\usepackage[hidelinks]{hyperref}

\newtheorem{theorem}{Theorem}
\newtheorem{lemma}{Lemma}
\newtheorem{remark}{Remark}

\newcommand{\F}{\mathcal F}
\newcommand{\Lip}{\operatorname{Lip}}

\title{A Direct Answer to Problem 3.12 of arXiv:2509.00722}
\author{Automatic Math Discovery Run \texttt{fa\_banach\_001}}
\date{\today}

\begin{document}
\maketitle

\section*{Source}

Richard J. Smith, \emph{Lipschitz-free spaces and Bossard's reduction argument},
arXiv:2509.00722.

The open problem appears as Problem 3.12 on page 11 of the source PDF.

\begin{center}
\includegraphics[width=\textwidth]{figures/open_problem_crop.png}
\end{center}

\section*{Problem Statement}

The paper asks:

\begin{quote}
Let \(M\) be a complete separable p1u metric space. Does there exist a
separable Banach space \(X\) such that \(\F(M) \hookrightarrow X^*\)?
\end{quote}

The notation section of the paper states that \(Y\hookrightarrow Z\), for
Banach spaces \(Y,Z\), means that \(Y\) is linearly isomorphic to a subspace of
\(Z\). The proof below answers the problem in this printed sense.

\section*{Result}

\begin{theorem}
Let \(M\) be a separable metric space. Then there exists a separable Banach
space \(X\) such that \(\F(M)\) embeds isometrically into \(X^*\). In fact one
may take \(X=\ell_1\).
\end{theorem}

\begin{proof}
First, \(\F(M)\) is separable whenever \(M\) is separable. Indeed, if
\(D\subset M\) is countable and dense and contains the base point, then the
linear span of \(\{\delta(x):x\in D\}\) is dense in \(\F(M)\).

We use the following standard fact. Let \(Y\) be a separable Banach space. There
is a countable norming subset of \(B_{Y^*}\). To see this, choose a dense
sequence \((y_n)\) in the unit sphere of \(Y\). For each \(n\), take
\(y_n^*\in S_{Y^*}\) with \(y_n^*(y_n)=1\). If \(y\in S_Y\) and
\(\varepsilon>0\), choose \(n\) with \(\|y-y_n\|<\varepsilon\). Then
\[
  |y_n^*(y)|\ge |y_n^*(y_n)|-|y_n^*(y-y_n)|>1-\varepsilon.
\]
It follows that
\[
  \|y\|=\sup_n |y_n^*(y)|\qquad (y\in Y).
\]

Define
\[
  T:Y\longrightarrow \ell_\infty,\qquad
  T(y)=(y_n^*(y))_{n=1}^\infty .
\]
The displayed equality says exactly that \(T\) is an isometric linear
embedding. Since \(\ell_\infty=(\ell_1)^*\) and \(\ell_1\) is separable, every
separable Banach space embeds isometrically into the dual of a separable
Banach space.

Apply this with \(Y=\F(M)\). This gives an isometric embedding
\[
  \F(M)\hookrightarrow \ell_\infty=(\ell_1)^*,
\]
so the required separable Banach space is \(X=\ell_1\).
\end{proof}

\section*{Verification}

The proof uses only two standard facts:

\begin{itemize}
\item Lipschitz-free spaces over separable metric spaces are separable.
\item Every separable Banach space has a countable norming subset of its dual
  unit ball, hence embeds isometrically into \(\ell_\infty=(\ell_1)^*\).
\end{itemize}

No computational verification is relevant here.

\section*{Limitations and Semantic Risk}

This packet solves Problem 3.12 as it is literally printed, using the paper's
own convention that \(\hookrightarrow\) denotes linear isomorphism to a
subspace.

The result does not address stronger variants that the author may have
intended but did not state, for example:

\begin{itemize}
\item requiring the dual space \(X^*\) itself to be norm-separable;
\item requiring \(\F(M)\) to be weak-star closed in \(X^*\);
\item requiring a complemented embedding;
\item asking whether \(\F(M)\) is itself isomorphic to a dual space.
\end{itemize}

Those variants are substantively different from the printed problem.

\section*{Verdict}

Under the exact notation and wording in arXiv:2509.00722, Problem 3.12 has an
affirmative answer for every separable metric space \(M\), not only for
complete separable p1u spaces.

\end{document}
